\chapter{Concluzii}

Aceasta lucrare a propus si aplicat un cadru metodologic pentru analiza
structurii de dependenta conditionala dintre sectoarele pietei financiare
americane. Abordarea combina modelarea econometrica prin modele vector
autoregresive cu estimarea grafica prin graphical lasso si analiza de retea
prin masuri de centralitate.

\section{Sinteza metodologica}

Structura metodologica a lucrarii poate fi rezumata in patru etape:

\begin{enumerate}
    \item \textbf{Preprocesarea datelor}: randamentele logaritmice zilnice ale
    celor 11 ETF-uri sectoriale au fost calculate pe esantionul analizat si
    organizate intr-o matrice de date multivariata.

    \item \textbf{Filtrarea VAR}: un model VAR(1), selectat prin criteriul BIC
    si validat prin teste de diagnostic, a eliminat dependenta temporala dintre
    randamente. Reziduurile astfel obtinute reprezinta socurile contemporane
    neanticipate ale fiecarui sector.

    \item \textbf{Estimarea retelei}: matricea de precizie a reziduurilor a
    fost estimata prin graphical lasso, utilizand un nivel de regularizare
    justificat prin analiza pe grila a parametrului $\rho$. Elementele nenule
    ale matricei de precizie definesc muchiile retelei financiare.

    \item \textbf{Analiza structurala}: masurile de centralitate (degree,
    strength, betweenness, eigenvector centrality) au permis cuantificarea
    importantei fiecarui sector in structura retelei. Analiza pe ferestre
    temporale a investigat stabilitatea in timp a acestei structuri.
\end{enumerate}

\section{Contributii principale}

Principalele contributii ale lucrarii pot fi sintetizate dupa cum urmeaza:

\begin{enumerate}
    \item A fost prezentat si argumentat riguros cadrul matematic al metodologiei,
    incluzand teoria modelelor VAR, proprietatile matricei de precizie in
    contextul grafurilor gaussiene si proprietatile de consistenta ale
    estimatorului GLasso.

    \item A fost demonstrat empiric ca structura pietei financiare americane
    nu este uniforma: anumite sectoare ocupa pozitii centrale cu grad de
    influenta mai ridicat, in timp ce altele sunt mai periferice.

    \item A fost ilustrat avantajul analizei bazate pe dependente conditionale
    fata de simpla analiza de corelatie: nu toate asocierile observate intre
    sectoare reflecta legaturi directe -- multe sunt mediate de factori comuni.

    \item Analiza dinamica pe ferestre temporale a aratat ca structura retelei
    nu trebuie privita ca perfect stabila in timp, ci ca susceptibila la
    modificari intre subperioade diferite ale esantionului analizat.
\end{enumerate}

\section{Concluzii statistice si matematice}

Din punct de vedere statistic, lucrarea ilustreaza utilitatea metodelor de
regularizare $\ell_1$ in contexte de inalta dimensionalitate. Proprietatea de
sparsitate indusa de penalizarea graphical lasso produce un estimator al
matricei de precizie care este simultan stabil numeric, consistent statistic
si interpretabil structural.

Din punct de vedere matematic, legatura dintre zerouri ale matricei de precizie
si independenta conditionala in distributia gaussiana multivariata ofera un
fundament teoretic riguros pentru interpretarea retelei financiare. Reteaua nu
este o constructie ad-hoc, ci un obiect cu proprietati matematice precise,
derivate din teoria grafurilor Markov gaussiene.

\section{Perspective}

Lucrarea poate fi extinsa in mai multe directii. Pe de o parte, esantionul ar
putea fi extins pentru a acoperi perioade mai lungi si mai multe cicluri economice.
Pe de alta parte, metodologia poate fi rafinata prin utilizarea unor modele
cu parametri variabili in timp sau prin incorporarea efectelor de volatilitate
conditionata. Compararea retelei de dependente conditionale cu reteaua de
cauzalitate Granger (construita in scriptul \texttt{04\_granger\_network.R})
reprezinta o alta directie valoroasa de investigare.

In ansamblu, lucrarea demonstreaza ca metodele matematice si statistice moderne
pot furniza o perspectiva structurata si riguroasa asupra organizarii interne
a pietei financiare, complementand si depasind informatia disponibila prin
masuri descriptive traditionale.